% GENERATED FILE. Edit canonical lessons, then run npm run generate:books.
% canonical-source-hash: fnv1a64:97516983641fd86c
% canonical-lessons: SA-C229-niyamah, SA-C229-dharmah, SA-C229-kartavyam, SA-C229-adhikarah, SA-C229-prasnapatram

\chapter{Rule, Righteousness, Obligation, Authority, Question paper}
\label{ch:sa-rule-righteousness-obligation-authority-question-paper}
\hlchaptermodality{229}

\begin{chapteropening}
\textbf{By the end of this chapter:} \emph{I can say a rule, righteousness, an obligation, authority and a question paper.}

Complete the last lesson of chapter 229: I can say a rule, righteousness, an obligation, authority and a question paper.
\end{chapteropening}

\section[\texorpdfstring{\sk{नियमः} (niyamaḥ)}{नियमः (niyamaḥ)}]{\sk{नियमः} (niyamaḥ) — a rule}

\label{lesson:SA-C229-niyamah}



\begin{quote}
Before the new one: say the Sanskrit for dance, then the Sanskrit for painting.
\end{quote}

\subsection*{You'll want to know: \sk{नियमः}}
\textbf{\sk{नियमः}} — \emph{niyamaḥ} — \textquotedblleft{}a rule\textquotedblright{}.

\begin{grammarlens}[title={What you've built}]
One more word for this chapter.
\end{grammarlens}

\subsection*{Your turn}
\begin{itemize}
\raggedright
  \item \emph{Say it:} \emph{niyamaḥ}
  \item \emph{Say it:} \emph{niyamaḥ}, once more
  \item \emph{Say it:} say \emph{niyamaḥ}
  \item \emph{Recall:} say the Sanskrit for dance, then the Sanskrit for painting, then say \emph{niyamaḥ} again
\end{itemize}

\begin{tcolorbox}[breakable,colback=teal!4,colframe=teal!35!black,title={Before you move on}]
What does \sk{नियमः} mean? (\textquotedblleft{}A rule\textquotedblright{}.) Say it once more.
\end{tcolorbox}
\section[\texorpdfstring{\sk{धर्मः} (dharmaḥ)}{धर्मः (dharmaḥ)}]{\sk{धर्मः} (dharmaḥ) — righteousness, duty}

\label{lesson:SA-C229-dharmah}



\begin{quote}
Before the new one: say the Sanskrit for painting, then the Sanskrit for a rule.
\end{quote}

\subsection*{You'll want to know: \sk{धर्मः}}
\textbf{\sk{धर्मः}} — \emph{dharmaḥ} — \textquotedblleft{}righteousness, duty\textquotedblright{}.

\begin{grammarlens}[title={What you've built}]
One more word for this chapter.
\end{grammarlens}

\subsection*{Your turn}
\begin{itemize}
\raggedright
  \item \emph{Say it:} \emph{dharmaḥ}
  \item \emph{Say it:} \emph{dharmaḥ}, once more
  \item \emph{Say it:} say \emph{dharmaḥ}
  \item \emph{Recall:} say the Sanskrit for painting, then the Sanskrit for a rule, then say \emph{dharmaḥ} again
\end{itemize}

\begin{tcolorbox}[breakable,colback=teal!4,colframe=teal!35!black,title={Before you move on}]
What does \sk{धर्मः} mean? (\textquotedblleft{}Righteousness, duty\textquotedblright{}.) Say it once more.
\end{tcolorbox}
\section[\texorpdfstring{\sk{कर्तव्यम्} (kartavyam)}{कर्तव्यम् (kartavyam)}]{\sk{कर्तव्यम्} (kartavyam) — an obligation, a duty}

\label{lesson:SA-C229-kartavyam}



\begin{quote}
Before the new one: say the Sanskrit for a rule, then the Sanskrit for righteousness.
\end{quote}

\subsection*{You'll want to know: \sk{कर्तव्यम्}}
\textbf{\sk{कर्तव्यम्}} — \emph{kartavyam} — \textquotedblleft{}an obligation, a duty\textquotedblright{}.

\begin{grammarlens}[title={What you've built}]
One more word for this chapter.
\end{grammarlens}

\subsection*{Your turn}
\begin{itemize}
\raggedright
  \item \emph{Say it:} \emph{kartavyam}
  \item \emph{Say it:} \emph{kartavyam}, once more
  \item \emph{Say it:} say \emph{kartavyam}
  \item \emph{Recall:} say the Sanskrit for a rule, then the Sanskrit for righteousness, then say \emph{kartavyam} again
\end{itemize}

\begin{tcolorbox}[breakable,colback=teal!4,colframe=teal!35!black,title={Before you move on}]
What does \sk{कर्तव्यम्} mean? (\textquotedblleft{}An obligation, a duty\textquotedblright{}.) Say it once more.
\end{tcolorbox}
\section[\texorpdfstring{\sk{अधिकारः} (adhikāraḥ)}{अधिकारः (adhikāraḥ)}]{\sk{अधिकारः} (adhikāraḥ) — authority, a right}

\label{lesson:SA-C229-adhikarah}



\begin{quote}
Before the new one: say the Sanskrit for righteousness, then the Sanskrit for an obligation.
\end{quote}

\subsection*{You'll want to know: \sk{अधिकारः}}
\textbf{\sk{अधिकारः}} — \emph{adhikāraḥ} — \textquotedblleft{}authority, a right\textquotedblright{}.

\begin{grammarlens}[title={What you've built}]
One more word for this chapter.
\end{grammarlens}

\subsection*{Your turn}
\begin{itemize}
\raggedright
  \item \emph{Say it:} \emph{adhikāraḥ}
  \item \emph{Say it:} \emph{adhikāraḥ}, once more
  \item \emph{Say it:} say \emph{adhikāraḥ}
  \item \emph{Recall:} say the Sanskrit for righteousness, then the Sanskrit for an obligation, then say \emph{adhikāraḥ} again
\end{itemize}

\begin{tcolorbox}[breakable,colback=teal!4,colframe=teal!35!black,title={Before you move on}]
What does \sk{अधिकारः} mean? (\textquotedblleft{}Authority, a right\textquotedblright{}.) Say it once more.
\end{tcolorbox}
\section[\texorpdfstring{\sk{प्रश्नपत्रम्} (praśnapatram)}{प्रश्नपत्रम् (praśnapatram)}]{\sk{प्रश्नपत्रम्} (praśnapatram) — a question paper}

\label{lesson:SA-C229-prasnapatram}



\begin{quote}
Before the new one: say the Sanskrit for an obligation, then the Sanskrit for authority.
\end{quote}

\subsection*{You'll want to know: \sk{प्रश्नपत्रम्}}
\textbf{\sk{प्रश्नपत्रम्}} — \emph{praśnapatram} — \textquotedblleft{}a question paper\textquotedblright{}.

\begin{grammarlens}[title={What you've built}]
One more word for this chapter.
\end{grammarlens}

\subsection*{Your turn}
\begin{itemize}
\raggedright
  \item \emph{Say it:} \emph{praśnapatram}
  \item \emph{Say it:} \emph{praśnapatram}, once more
  \item \emph{Say it:} say \emph{praśnapatram}
  \item \emph{Recall:} say the Sanskrit for an obligation, then the Sanskrit for authority, then say \emph{praśnapatram} again
\end{itemize}

\begin{tcolorbox}[breakable,colback=teal!4,colframe=teal!35!black,title={Before you move on}]
What does \sk{प्रश्नपत्रम्} mean? (\textquotedblleft{}A question paper\textquotedblright{}.) Say it once more.
\end{tcolorbox}
